\chapter{Introduzione}
\section{Contesto e funzione di una incubatrice neonatale}
\subsection{Esigenze dei neonati prematuri}
Con il termine di ``neonati prematuri'' si intendono tutti i bambini nati prima della 37\textsuperscript{a} settimana di
gestazione. A causa del loro mancato sviluppo, non dispongono di un sistema di termoregolazione adeguato a mantenere la
temperatura corporea ideale compresa tra i 36,5 °C e i 37,5 °C. Già quando scende al di sotto dei 36°C, il neonato entra
in una fase di ipotermia e se tale condizione si protrae nel tempo, si possono verificare gravi conseguenze sulla salute
del neonato, fino a portarlo alla morte.

Si identifica la temperatura critica come la temperatura a cui un essere umano non è più in grado di regolare la propria
temperatura corporea. Per un adulto si aggira attorno agli 0°C mentre per un neonato prematuro è compresa tra i 20°C e
i 23°C. Dato che questo valore è particolarmente elevato, casi di ipotermia sono molto frequenti nei neonati prematuri
in qualsiasi parte del mondo, anche nei paesi tropicali.

Per mantenere il neonato in condizioni di salute ottimali, risulta, quindi, necessario adottare misure per limitare
la dispersione di calore. Queste possono essere semplici azioni fattibili da chiunque, come il mantenimento del neonato
asciutto ed avvolto da indumenti caldi, oppure il metodo canguro-madre che consiste nel contatto pelle a pelle tra il
neonato e la madre. In alternativa esistono soluzioni più avanzate come l'uso di incubatrici neonatali o culle riscaldate,
in grado di mantenere una temperatura costante e confortevole per il neonato. \cite{world1993thermal}

\subsection{Storia delle incubatrici neonatali}
La prima idea di incubatrice neonatale venne proposta nel 1857 da Jean-Louis-Paul Denucé (1824 - 1889), un medico belga,
che, per mantenere in vita (purtroppo invano) un neonato prematuro, per primo costruì e descrisse formalmente una
culla di metallo a doppia parete senza coperchio superiore, in cui, nella cavità tra le due pareti, veniva versata acqua
calda. Tale struttura di incubatrice, però, non ebbe grande diffusione.

Nel 1881, un medico francese, Stéphane Tarnier (1828 - 1897), durante una visita allo Zoo di Parigi notò un'incubatrice
chiusa usata per le uova di uccelli e subito ne intuì le potenzialità per i neonati prematuri. Fece costruire un
modello simile con pareti in legno, un compartimento inferiore per contenere l'acqua calda e introdusse per la prima volta
il coperchio in vetro rendendola chiusa e più efficace. La struttura di Tarnier venne poi resa più semplice ed economica
da Alfred Auvard (1855-1940) negli anni successivi, rendendola più accessibile agli ospedali e favorendone la rapida
diffusione. Grazie a questo modello, Tarner riuscì a far aumentare la probabilità di sopravvivenza dei neonati di peso
inferiore ai 2kg dal 35\% al 62\%. \cite{DunnF137}

Poco più tardi, il medico americano Joseph Bolivar DeLee (1869 - 1942), nel 1899 durante il suo servizio presso il Chicago
Lying-In Hospital, notò la necessità di dover fornire assistenza tempestiva e sul posto ai neonati prematuri, che nascevano
a casa. Decise, quindi, di modificare il modello di Auvard per creare la prima incubatrice portatile, facilmente trasportabile
dall'ospedale al domicilio del neonato.

Ad oggi le incubatrici neonatali sono molto diffuse ed utilizzate. Sono costruite da una struttura trasparente di policarbonato
con delle aperture per l'accesso rapido al neonato. Dispongono di sistemi di riscaldamento, ventilazione e umidificazione,
oltre che a numerosi sensori per controllare lo stato dell'incubatrice e del neonato e a sistemi di allarme per segnalare
eventuali malfunzionamenti o criticità.

\section{Obiettivi della tesi}
L'obiettivo principale di questa tesi consiste nel rendere completamente funzionante il prototipo di incubatrice neonatale
costruito da Vanni Bincoletto nella sua tesi magistrale in bioingegneria \cite{bincoletto2025} e sviluppare un sistema di
gestione della temperatura.


\begin{itemize}
	\item sviluppo multidisciplinare e trasversale (hardware, software, controllore)
	\item mettere in pratica le conoscenze teoriche viste a lezione
	\item verificare l'efficacia dei controllori pid per il controllo della temperatura
\end{itemize}

\section{Struttura dell'elaborato}
Il seguente elaborato è organizzato in quattro capitoli principali, oltre al primo capitolo di introduzione. In ognuno di
essi viene trattato un argomento specifico corrispondente ad una fase dello sviluppo del prototipo di incubatrice neonatale.

Nel secondo capitolo vengono descritti i vari componenti elettronici del prototipo come il microcontrollore, i sensori
ed i vari attuatori. Vengono inoltre presentate le connessioni realizzate tra di essi per garantire il corretto funzionamento
di ogni componente del sistema.

Nel terzo capitolo si affronta lo sviluppo del firmware per il microcontrollore. Vengono enunciati i requisiti che hanno
guidato le scelte implementative e si analizzano i vari moduli in cui è stato suddiviso il codice.

Nel quarto capitolo si tratta del sistema di controllo usato per la gestione della temperatura. Viene prima descritto un
semplice controllore ad isteresi e successivamente si passa ad un controllore PID. Si conclude con un confronto tra i due
metodi di controllo, evidenziando i vantaggi, gli svantaggi ed i limiti di ognuno.

Nel quinto capitolo si discutono in maniera critica i risultati raggiunti rispetto agli obiettivi prefissati all'inizio
e si propongono possibili sviluppi futuri per migliorare ulteriormente il prototipo di incubatrice neonatale.

A concludere è presente un'appendice in cui sono riportati la dimostrazione matematica di una formula usata nel controllore
PID e il codice sorgente completo del firmware.

In questa tesi verrà analizzata esclusivamente la gestione del sistema di riscaldamento. Poiché, tuttavia, il prototipo
di incubatrice neonatale è dotato anche di un umidificatore, è stata prevista la possibilità di implementare in futuro
un sistema di controllo dell’umidità. A tal fine, il possibile utilizzo dell’umidificatore è stato considerato sia nella
progettazione della circuiteria hardware sia nello sviluppo del software, senza implementare un valido sistema di controllo
dedicato alla gestione dell’umidità.
